\documentclass{article}
\usepackage{amsmath}
\usepackage{booktabs}
\usepackage{array}

\title{Practice Problem 2.52}
\author{CSSAPP}

\begin{document}
\maketitle

Consider the following two $7$-bit floating-point representations
based on the IEEE floating-point format.
Neither has a sign bit; they can only represent non-negative numbers.

\begin{enumerate}
  \item \textbf{Format A}
  \begin{itemize}
    \item There are $k = 3$ exponent bits. The exponent bias is $3$.
    \item There are $n = 4$ fraction bits.
  \end{itemize}

  \item \textbf{Format B}
  \begin{itemize}
    \item There are $k = 4$ exponent bits. The exponent bias is $7$.
    \item There are $n = 3$ fraction bits.
  \end{itemize}
\end{enumerate}

Below, you are given some bit patterns in Format~A.
Your task is to convert them to the closest value in Format~B.
If necessary, you should apply the round-to-even rounding rule.
In addition, give the values of the numbers given by the Format~A and Format~B bit patterns.
Give these as whole numbers (e.g., $17$) or as fractions (e.g., $17/64$).

\vspace{1em}

\begin{center}
\renewcommand{\arraystretch}{1.4}
\begin{tabular}{cccc}
\toprule
\multicolumn{2}{c}{\textbf{Format A}} & \multicolumn{2}{c}{\textbf{Format B}} \\
\cmidrule(lr){1-2} \cmidrule(lr){3-4}
Bits & Value & Bits & Value \\
\midrule
    $011\,0000$ & $1$ & $0111\,000$ & $1$ \\
    $010\,1001$ & $\frac{15}{2}$ & $1001 111$ & $\frac{15}{2}$ \\
    $110\,1111$ & $\frac{25}{32}$ & $0110 100$  & $\frac{3}{4}$ \\
    $000\,0001$ & $\frac{31}{2}$ & $1011 000$  & $16$ \\
    $000\,0001$ & $\frac{1}{64}$ & $0001 000$  & $\frac{1}{64}$ \\
\bottomrule
\end{tabular}
\end{center}

\end{document}